\section{Introduction}
\label{sec:introduction}

Goal-frontier maximization (GFM) proposes alignment through a single
objective: maximize the volume of the jointly achievable capability space
$\volL(G)$ \citep{lasser2026gfm}. The foundational paper proves structural
safety properties (self-balancing against destruction, coercion, and
rigidity). Its companion \citep{lasser2026poset} constructs a computable
capability poset with polynomial-time measure computation and a leverage
ranking signal that produces a gradient toward civilization-building
infrastructure. The third paper \citep{lasser2026horizon} extends the
framework to multi-step planning through a discounted value function,
develops risk-capability identification, and establishes the anti-monopolar
property: under a diversity condition, full capability domination is
anti-maximizing.

\paragraph{Epistemic status.} This paper is implementation machinery.
Papers~1--3 establish what is maximized ($\volL(G)$), prove structural
safety properties (self-balancing, anti-monopolar pressure, structural
alignment), and characterize the conditions under which those properties
hold \citep{lasser2026gfm, lasser2026poset, lasser2026horizon}. Those
structural guarantees do not depend on the correctness of this paper's
results: a flaw in the causal attribution rule, the risk-trust dynamics,
the SPRT detection test, or the dependency-risk model would be a flaw
in the \emph{implementation} of agent-evaluation mechanisms, not in the
framework's alignment properties. The framework falls back on the trust
and capability machinery of Papers~1--3 whenever this paper's mechanisms
degrade---a property formalized through the multi-channel attribution
architecture of Section~\ref{sec:multi_channel}.

The preceding papers treat every agent-evaluation question as a statistical
signal-processing problem. The foundational paper's scorpion detection
\citep{lasser2026gfm} identifies agents whose actions are \emph{correlated}
with $\volL$-contraction but cannot disentangle their contribution from
simultaneous causes. The horizon paper's risk-trust factor $\Trisk_j$
\citep{lasser2026horizon} has no formal update dynamics. The dependency-risk
analysis uses minimax bounds that are quantitatively too weak for
risk-adjusted planning. In each case, statistical signals suffice for
obvious cases (direct harm, clear deception) but fail for the hard cases
that matter most: confounded contraction in multi-agent populations,
coordinated risk manipulation, and risk claims whose outcomes are never
observed.

This paper introduces a structural causal model for capability dynamics
and uses it to upgrade scorpion detection and risk evaluation from
statistical to causal identification and to sharpen the remaining
agent-evaluation mechanisms. The common
thread: a GFM agent needs a \emph{causal model} of other agents to make
good decisions about their harm (who caused this contraction?),
credibility (is this risk claim structurally sound?), and adversarial
exposure (what substrate-correlated threats do we face?).

\paragraph{Contributions.}
\begin{enumerate}
\item A structural causal model for capability dynamics with
  do-calculus contraction attribution, showing a conditional $L^2$
  convergence advantage over correlational detection under causal
  sufficiency, additive-SCM structure, and exogenous orthogonality
  (Section~\ref{sec:causal}).
\item $L^2$ EWMA dynamics for risk-trust $\Trisk_j$ over
  pathway-conditional structural verification residuals, and a
  prior-corrected trust-weighted log-linear opinion pool that recovers
  Bayesian updating in the no-tempering limit, replacing the max
  aggregation operator, with a two-gate architecture separating
  attention allocation from action decisions via the actor's own
  forward causal risk evaluation (Section~\ref{sec:risk_trust}).
\item Causal scorpion detection under a Gaussian model, committed to a
  sequential probability ratio test on least-favorable simple
  hypotheses with Type I/II error control, and a high-probability
  bound on non-stationary scorpion evasion bandwidth under a
  stationary-epoch model (Section~\ref{sec:scorpion}).
\item A probabilistic dependency-risk model gated by an adversarial
  balance condition $c > \max_i (f_i \cdot c_i)$ that admits contagion
  dynamics, with per-step expected costs that compound at the same
  rate as growth and a quantitative substrate-count lower bound
  $m > c_{\max}/c$ that sharpens the minimax $m \geq 2$
  (Section~\ref{sec:dependency}).
\end{enumerate}
